%! AP Statistics — Unit 4, Lesson 4.3 — Group Activity KEY
\documentclass[10pt]{article}
\usepackage{apstats-article}
\usepackage{apstats-key}

\begin{document}

\pageheader{Unit 4, Lesson 4.3}{Group Activity --- KEY}
\namepartnerperiod

\vspace{0.1in}

\begin{scenariobox}[Context]{sky}
\small
A standard deck of 52 playing cards has 4 suits (hearts, diamonds, clubs, spades),
13 cards per suit (Ace, 2--10, Jack, Queen, King). Hearts and diamonds are red;
clubs and spades are black. Face cards = Jack, Queen, King (3 per suit, 12 total).
\end{scenariobox}

\vspace{0.1in}

% ── Tier R ────────────────────────────────────────────────────────────────────
\begin{tcolorbox}[colback=goldbg, colframe=goldacc, boxrule=1pt, arc=2mm,
    left=4mm, right=4mm, top=3mm, bottom=3mm,
    title={\bfseries Tier R \enspace Remediate --- Building Sample Spaces}]
\small

\begin{enumerate}[leftmargin=1.5em, itemsep=8pt]
  \item \ans{52 total outcomes.}

  \item
        \renewcommand{\arraystretch}{1.8}
        \begin{tabularx}{\linewidth}{@{} l X @{}}
        Event & $P(E) =$ \\
        \hline
        Drawing a heart & \ans{$13/52 = 1/4 = 0.25$} \\
        Drawing a red card & \ans{$26/52 = 1/2 = 0.5$} \\
        Drawing a face card & \ans{$12/52 = 3/13 \approx 0.231$} \\
        Drawing the Ace of Spades & \ans{$1/52 \approx 0.019$} \\
        \end{tabularx}

  \item
        \begin{enumerate}[label=(\alph*), itemsep=6pt]
          \item \ans{$P(\text{not a heart}) = 1 - 1/4 = 3/4 = 0.75$}
          \item \ans{$P(\text{not a face card}) = 1 - 3/13 = 10/13 \approx 0.769$}
        \end{enumerate}

  \item \ans{$P = 0$: drawing a card with 7 suits (impossible). $P = 1$: drawing any card at all (certain, since every draw produces a card). Both are possible to state but the only realistic example from this deck is $P(\text{draw a valid card}) = 52/52 = 1$.}
\end{enumerate}

\end{tcolorbox}

\vspace{0.1in}

% ── Tier A ────────────────────────────────────────────────────────────────────
\begin{tcolorbox}[colback=greenbg, colframe=greenacc, boxrule=1pt, arc=2mm,
    left=4mm, right=4mm, top=3mm, bottom=3mm,
    title={\bfseries Tier A \enspace Approaching Proficiency --- Probability and Interpretation}]
\small

\begin{enumerate}[leftmargin=1.5em, itemsep=8pt]
  \item
        \renewcommand{\arraystretch}{1.8}
        \begin{tabularx}{\linewidth}{@{} l p{2.5cm} X @{}}
        Event & $P(E)$ & Long-run interpretation \\
        \hline
        Drawing a club & \ans{$13/52 = 1/4$} & \ans{If a card is drawn and replaced many times, about 25\% of draws will be clubs.} \\
        Drawing a King & \ans{$4/52 = 1/13$} & \ans{In the long run, about 1 in every 13 draws will be a King.} \\
        Drawing a red face card & \ans{$6/52 = 3/26$} & \ans{About 3 out of every 26 draws will be a red face card in the long run.} \\
        \end{tabularx}

  \item \ans{The error is confusing the \emph{outcome} of one specific draw with the \emph{probability} of the event. The probability is a long-run measure, not a statement about what just happened. One draw of a heart does not change the theoretical probability $P(\text{heart}) = 1/4$.}

  \item
        \begin{enumerate}[label=(\alph*), itemsep=6pt]
          \item \ans{$P(\text{not a King}) = 1 - 4/52 = 48/52 = 12/13 \approx 0.923$}
          \item \ans{$P(\text{not a red face card}) = 1 - 6/52 = 46/52 = 23/26 \approx 0.885$}
        \end{enumerate}

  \item \begin{teachernote}Tree: Color branch first (Red/Black), then suit branch (H, D for Red; C, S for Black). Each of the 4 equally-likely suit branches represents 1/4 of the deck. $P(\text{red}) = P(\text{heart}) + P(\text{diamond}) = 1/4 + 1/4 = 1/2$.\end{teachernote}
\end{enumerate}

\end{tcolorbox}

\vspace{0.1in}

% ── Tier E ────────────────────────────────────────────────────────────────────
\begin{tcolorbox}[colback=sky, colframe=navy, boxrule=1pt, arc=2mm,
    left=4mm, right=4mm, top=3mm, bottom=3mm,
    title={\bfseries Tier E \enspace Extension --- Multi-Stage Sample Spaces}]
\small

\begin{enumerate}[leftmargin=1.5em, itemsep=8pt]
  \item \ans{$10 \times 10 = 100$ outcomes.}
  \item \ans{Spin 1 and Spin 2 each land on sections 1--5 (red). Number of outcomes: $5 \times 5 = 25$.}
  \item \ans{$P(\text{both red}) = 25/100 = 1/4 = 0.25$.}
  \item \ans{$P(\text{no blue on either spin}) = P(\text{not blue on spin 1}) \times P(\text{not blue on spin 2})$. Sections NOT blue: 7 out of 10 each spin. $P(\text{no blue}) = (7/10)^2 = 49/100$. So $P(\text{at least one blue}) = 1 - 49/100 = 51/100 = 0.51$.}
  \item \ans{If the spinner is spun twice many times, about 51\% of the time at least one of the two spins will land on a blue section.}
  \item \ans{No --- the two spins are independent. Knowing the first spin was red provides no information about the second spin's outcome. The probability that spin 2 is red remains $5/10 = 0.5$ regardless of what spin 1 showed.}
\end{enumerate}

\end{tcolorbox}

\end{document}
